\section{Audited common-factorial theorem for the prospective 26
pairs}

\subsection{1. Scope and claim status}

This note supplies the proof requested by the claim-neutral selector in
\texttt{src/\texttt{\seqsplit{prospective\_no\_promotion\_26.py}}}. It proves two local
extensions and then composes them on a fixed irreducible class:

\begin{enumerate}
\def\labelenumi{\arabic{enumi}.}
\tightlist
\item
  the audited arbitrary-orientation one-active graph/kernel theorem
  applies to the selector's 30 one-active incidences, with the same
  arbitrary fixed linear correction used below;
\item
  the curvature-safe reversible two-node all-active estimate lifts to
  the \textbf{discrete} rate-adjusted factorial and to its fourth power;
  and
\item
  one common potential closes the one-active and all-active descriptions
  by the existing all-species reflected-debt gluing theorem.
\end{enumerate}

The resulting pair statement has passed an independent proof audit. The
finite scope certificate is
\texttt{\seqsplit{src/\texttt{\seqsplit{prospective\_26\_candidate\_theorem.py}}}}, and its focused
tests are in
\texttt{\seqsplit{tests/\texttt{\seqsplit{test\_prospective\_26\_candidate\_theorem.py}}}}. Its two
scoped analytic flags and its 26-pair recurrence flag are certified. The
global T3-2 flag remains false.

There is one correction to the older selector note. Proposition 5.2 of
\texttt{three\_\texttt{\seqsplit{active\_shell\_gluing\_gate.md}}} is literally written for
the continuous entropy

\[
 U_\theta(x)=\sum_i\left[x_i\left(\log{x_i\over\theta_i}-1\right)
                         +\theta_i\right],
\]

not for \(\sum_i\log(x_i!)\). The two functions do not differ by an
affine term. Section 4 below proves the missing discrete comparison and
does not identify the two potentials by assertion.

\subsection{2. Exact selector and the common
correction}

The selector contains 26 positive-invariant ordered support pairs, with
pair fingerprint

\begin{CodeBlock}
393474671be0bf095868e66cbcbf3164d941b99191517f172a41f157e20b21af
\end{CodeBlock}

Their 124 affine-feasible failed incidences split as

\[
 30\text{ one-active},\qquad 0\text{ two-active},\qquad
 94\text{ all-active}.                                      \tag{2.1}
\]

The selector is disjoint from every prior-certified branch and lies
inside the authoritative claim-neutral prospective 795. The frozen
arithmetic is

{\small\def\LTcaptype{none} % do not increment counter
\begin{longtable}[]{@{}lrrr@{}}
\toprule\noalign{}
claim-neutral set & positive & signed & total \\
\midrule\noalign{}
\endhead
\bottomrule\noalign{}
\endlastfoot
prospective parent & 759 & 36 & 795 \\
selected candidate, prior-certified overlap zero & 26 & 0 & 26 \\
prospective remainder after candidate & 733 & 36 & 769 \\
\end{longtable}
}

The 769-pair remainder fingerprint is

\begin{CodeBlock}
4e645eca1ba23849680f2f983e1fb8c9465001c5b0e8c0090a31a339bf18ec06
\end{CodeBlock}

Before the independent verdict this was prospective arithmetic only. The
audit has now promoted exactly these 26 pairs, with zero overlap, and
changes the authoritative remainder from \((759,36)\) to \((733,36)\).

Fix one pair, arbitrary strongly connected orientations of both linkage
supports, arbitrary positive rates, and one closed irreducible
population class \(\Gamma\). The whole-top linkage in every failed
all-active descriptor of this pair is one fixed support \(T=\{y,z\}\).
It is either

\[
 \{2A,B+C\}\quad\hbox{or}\quad\{A+C,B+C\}.           \tag{2.2}
\]

A strongly connected orientation on two vertices contains both
directions. After aggregating any parallel labels, write their positive
rates as \(\kappa_{yz}\) and \(\kappa_{zy}\), and put \(d=z-y\). Choose
\(\theta\in(0,\infty)^3\) so that

\[
 \kappa_{yz}\theta^y=\kappa_{zy}\theta^z,
 \qquad
 \ell_i=-\log\theta_i.                              \tag{2.3}
\]

Only the scalar equation

\[
 \ell\mathbin\cdot d
   =-\log(\kappa_{yz}/\kappa_{zy})                   \tag{2.4}
\]

is prescribed, so a solution always exists. Because \(T\) is fixed for
the pair, the same \(\ell\) works on all 94-menu cones belonging to that
pair. Choose \(K\) large enough that

\[
 F_\ell(x)=K+\sum_{i=1}^3\log(x_i!)+\ell\mathbin\cdot x,
 \qquad G=1+F_\ell\ge1,
 \qquad W=G^4                                           \tag{2.5}
\]

on \(\mathbb N_0^3\). Such a \(K\) exists because the factorial term
dominates every fixed linear term. The function \(W\) is proper and

\[
 G(x)\asymp R\log R,qquad R=2+\lVert x\rVert_1,      \tag{2.6}
\]

along every divergent sequence.

\subsection{3. The 30 one-active rows}

\subsubsection{3.1 Exact support audit}

Relabel the active species as \(X=C\) and the inactive species as
\(A,B\), as in the audited one-active graph theorem. The 30 physical
rows collapse to 15 exact normalized support profiles. Every row has
capped inactive state \((A,B)=(0,0)\), and both linkages are mixed.
Their exact partition is

{\small\def\LTcaptype{none} % do not increment counter
\begin{longtable}[]{@{}lr@{}}
\toprule\noalign{}
arbitrary-orientation route & incidences \\
\midrule\noalign{}
\endhead
\bottomrule\noalign{}
\endlastfoot
a mixed linkage contains the pure active source \(X\) & 20 \\
Family III, resistance-zero service from the origin & 8 \\
Family III, frozen origin/no positive-debt history & 2 \\
\end{longtable}
}

The full 30-row and 15-profile hashes are

\begin{CodeBlock}
one-active rows:     091dc59dc83b6445d36135a2d6b1e6c1bfb9565db88328fc1725c39f78bbc7e4
normalized profiles: 36b6eb1c1beadc939681fc7515e2adf3dfb9e16d116da4b8d2bd73ac77aabe12
\end{CodeBlock}

The certificate checks the following support facts row by row.

\begin{itemize}
\tightlist
\item
  In each of the 20 direct rows, one mixed linkage has stripped top
  source \(0\), hence physical source \(X\), and also has a lower
  complex.
\item
  In each of the ten remaining rows the two stripped top sources are
  exactly \(A\) and \(B\), one in each mixed linkage. The no-fast set is
  therefore the origin.
\item
  In the eight service rows, a zero-containing linkage has another lower
  vertex. In the other two rows neither lower support contains \(0\).
\end{itemize}

These are exactly the hypotheses used in Sections 3 and 6 of the audited
arbitrary-orientation graph proof. They do not depend on a
Hamilton-cycle orientation or a bounded CEGAR box.

\subsubsection{Lemma 3.1 (scope extension of the graph
ordering)}

For every strong orientation and every positive rate vector, every
historically consistent positive-debt row among these 30 incidences has

\[
 m_-=0<m_+,                                             \tag{3.1}
\]

where resistance counts a lower-source firing while an active-degree-one
source is enabled. The two remaining rows are frozen at the displayed
origin and admit no history carrying positive reflected \(X\)-debt.

\paragraph{Proof}

In a direct row, positive old debt implies \(X\ge1\), so the source
\(X\) is enabled. In its strong mixed linkage take a simple directed
complex path from \(X\) to a lower complex. Before the first lower
target every source has active degree one. The top-to-top edges have
active reward zero and the first top-to-lower edge has reward \(-1\).
Consecutive edges are physically enabled because the target complex of
one firing supplies the source of the next. This is a resistance-zero
strict decrease.

In a Family-III service row, the only no-fast base is the origin. Follow
a simple path in the zero-containing strong linkage from \(0\) to a
nonzero lower vertex. If the path enters an active complex first, follow
its first exit; the surviving nonzero cofactor enables the other active
source or a second exit. In either case the path supplies one more exit
than entry and strictly lowers the relative active level without firing
a lower source against an enabled fast clock. Hence again \(m_-=0\).

A primitive positive return to the origin cannot have resistance zero.
Its only free launch occurs while no fast source is enabled. Thereafter,
in a zero-resistance word, every reaction is active-sourced until the
origin is reached. Each entry contributes \(+1\), each exit contributes
\(-1\), and return of the inactive population requires at least as many
exits as free entries. Its active reward is therefore nonpositive. Thus
\(m_+\ge1\).

Finally, in either no-history row the origin has no enabled source: the
two top sources require \(A\) and \(B\), and every lower source is
nonzero. Moreover no linkage contains target \(0\), so this origin
cannot be reached from a nonorigin state. It is an absorbing initial
singleton and cannot carry positive reflected debt. This proves the
lemma. \(\square\)

\subsubsection{\texorpdfstring{Lemma 3.2 (aggregate kernel and arbitrary
fixed
\(\ell\))}{Lemma 3.2 (aggregate kernel and arbitrary fixed \textbackslash ell)}}

For every nonfrozen row of Lemma 3.1, the complete physical one-active
block of the audited fourth-power kernel theorem satisfies, uniformly
over the finite base marks of the fixed class,

\[
 p_D(n)\ge a,
 \qquad p_U(n)\le {b\over n},
 \qquad \mathbb E\sigma_n\le T,                     \tag{3.2}
\]

for all sufficiently large active population \(n\). Every fixed
endpoint-polynomial moment is bounded in the sense of the kernel
theorem, and its stopped three-interruption moving-boundary remainder is
\(O((L_n^2/n)^3)\) with \(L_n=n^{1/8}\). The conclusion remains true for
the common \(W\) in (2.5), with the arbitrary fixed \(\ell\) selected in
(2.3).

\paragraph{Proof}

The analytic proof of the resistance-to-kernel proposition uses only the
support properties just checked, binary molecularity, strong
connectivity, and fixed positive rates.

In a direct row, divide time by the active population. Before service,
the stripped active-source phase is an immigration--conversion--death
process on complexes contained in \(\{0,A,B\}\), killed on its first
active exit. Every transient top component has a directed route to the
lower part because the original mixed linkage is strong. The exponential
killed-Green estimate and its polynomially size-biased versions
therefore apply exactly as in the audited direct-block proof. The simple
path in Lemma 3.1 has a strictly positive product of rate fractions,
which supplies the order-one down coefficient.

In a Family-III service row the regenerative no-fast set is the single
origin. Between visits to it, the same killed unimolecular Green
estimate applies. The resistance-zero witness supplies the positive
order-one down coefficient. A positive origin return requires at least
one suppressed lower clock by Lemma 3.1, so the ordered-compensation
expansion makes its aggregate probability \(O(n^{-1})\), after all
neutral loops have been summed. Binary molecularity gives the same
carrier estimate and the same endpoint-weighted three-interruption bound
as in the audited proof. There is no unlisted countable no-fast phase in
these rows.

This proves (3.2), including duration and endpoint moments, rather than
only the probability of one selected word. The repeated-base
fourth-power theorem now applies with \(m=0\).

It remains to check \(\ell\). A bounded reaction changes
\(\ell\mathbin\cdot x\) by the fixed number
\(\ell\mathbin\cdot\zeta_r\). Thus it changes neither resistance, the
killed generator, nor any power of \(n\) in the compensation expansion.
At a down endpoint,

\[
 F_\ell(X_{\tau})-F_\ell(X_0)=-\log n+O(1+Z),       \tag{3.3}
\]

where \(Z\) is the inactive endpoint cost; at an up endpoint the
analogous positive bound differs from the \(\ell=0\) bound by at most a
fixed multiple of the endpoint population. The already proved
exponential/factorial endpoint moments absorb that linear term. At the
moving boundary the term is \(O(L_n+J)\), not merely \(O(L_n)\): nested
active entries can make the active overshoot as large as a constant
multiple of the counted number \(J\) of suppressed interruptions. The
audited boundary estimate is endpoint-weighted and controls every fixed
\(J\)-moment, in particular an order strictly above eight, on the
three-interruption remainder. Applying that estimate first to the
leading linear endpoint term and then with moment order \(q>8\) to the
fourth-power tail gives the same
\(O(n^{2+7/8}(\log n)^4)=o(n^3(\log n)^4)\) completed-episode boundary
cost as for \(\ell=0\). Hence the fourth-power endpoint calculation and
the promotion-access estimate are unchanged. \(\square\)

\subsection{4. Discrete all-active fourth-power
lift}

Fix one of the 94 all-active rows and a tier sequence \(x_n\) realizing
it. All three coordinates diverge. Hence, after deleting a finite
prefix, both top sources are enabled, both neighboring states
\(x_n\pm(z-y)\) lie in the population lattice whenever they occur below,
and all factorial quotients are literal. At boundary states the usual
convention is \(x^{\underline y}=0\) when \(y\nleq x\); no boundary
state remains in this all-active tail. Let \(T=\{y,z\}\), \(d=z-y\), and
write

\[
 A(x)=\kappa_{yz}x^{\underline y},\qquad
 B(x)=\kappa_{zy}x^{\underline z},\qquad
 h(x)=\log{A(x)\over B(x)}.                         \tag{4.1}
\]

The two top source monomials occupy one D-tier, so along a further tier
subsequence

\[
 0<c\le A/B\le C<\infty.                            \tag{4.2}
\]

Let \(R\) now denote the other linkage, let \(\beta_n\) be its maximal
source propensity, and put

\[
 \mathcal D_n=(A_n-B_n)\log(A_n/B_n)\ge0.           \tag{4.3}
\]

\subsubsection{Lemma 4.1 (discrete reversible
decomposition)}

For the top linkage,

\[
 \mathcal L_TF_\ell(x_n)
 \le-\mathcal D_n+C\beta_n,                         \tag{4.4}
\]

and, for \(k=2,3,4\),

\[
 \sum_{r\in T}\lambda_r(x_n)|\Delta_rF_\ell(x_n)|^k
 \le C(\mathcal D_n+\beta_n).                      \tag{4.5}
\]

\paragraph{Proof}

For the forward jump, the factorial identity is exact:

\[
 \prod_i{(x_i+d_i)!\over x_i!}
 ={(x+d)^{\underline z}\over x^{\underline y}}.   \tag{4.6}
\]

Together with (2.4), it gives

\[
 e^{F_\ell(x+d)-F_\ell(x)}
 ={\kappa_{zy}(x+d)^{\underline z}
    \over \kappa_{yz}x^{\underline y}},            \tag{4.7}
\]

the exact discrete endpoint detailed-balance ratio. To compare both top
clocks at the same state, bounded falling-factorial expansion gives

\[
 \Delta_{y z}F_\ell=-h+\varepsilon_+,qquad
 \Delta_{z y}F_\ell=h+\varepsilon_-,               \tag{4.8}
\]

where, with \(I=\{i:d_i\ne0\}\),

\[
 |\varepsilon_+|+|\varepsilon_-|
 \le C s_n,qquad s_n=\sum_{i\in I}{1\over x_{n,i}}. \tag{4.9}
\]

Indeed each logarithm in the quotient of two shifted factorials is a sum
of finitely many terms \(\log(1+j/x_i)\).

Consequently

\[
 \mathcal L_TF_\ell
 =-(A-B)h+A\varepsilon_++B\varepsilon_-.
                                                               \tag{4.10}
\]

The curvature-safe premise is exactly

\[
 {A_n+B_n\over x_{n,i}}\le C\beta_n
 \qquad(i\in I).                                    \tag{4.11}
\]

To see this, choose a top endpoint \(e\) with \(e_i>0\). Top flatness
makes \(A+B\asymp x_n^e\), while \(x_n^e/x_{n,i}=x_n^{e-\mathbf e_i}\);
the certificate checks that this curvature cofactor lies at or below the
maximal \(R\)-source D-tier in all 94 rows. Equations (4.9)--(4.11)
prove (4.4).

On the compact ratio range (4.2),

\[
 \mathcal D_n\asymp(A_n+B_n)h_n^2.                 \tag{4.12}
\]

Equations (4.8)--(4.9), boundedness of \(h_n\), and
\((A+B)s_n^2\le C\beta_n\) therefore yield, for \(2\le k\le4\),

\[
 (A+B)(|\Delta_{yz}F_\ell|^k+|\Delta_{zy}F_\ell|^k)
 \le C\{(A+B)h^2+(A+B)s_n^2\}
 \le C(\mathcal D_n+\beta_n),
\]

which is (4.5). \(\square\)

\subsubsection{Lemma 4.2 (lower-linkage drift and jump
moments)}

There is a sequence \(a_n\to\infty\) such that

\[
 \mathcal L_RF_\ell(x_n)\le-\beta_na_n,             \tag{4.13}
\]

and, for \(k=2,3,4\),

\[
 \sum_{r\in R}\lambda_r(x_n)|\Delta_rF_\ell(x_n)|^k
 \le C\beta_n(1+\log^k R_n),
 \qquad R_n=2+\lVert x_n\rVert_1.                  \tag{4.14}
\]

\paragraph{Proof}

The lower-linkage part of Proposition 5.2 proves
\(\mathcal L_RU_\theta/\beta_n\to-\infty\). Here

\[
 F_\ell-U_\theta
 =C_0+\sum_i\{\log(x_i!)-x_i\log x_i+x_i\}.        \tag{4.15}
\]

For every bounded reaction vector, the jump of the right side is
\(O(\sum_i x_i^{-1})=o(1)\) on an all-active sequence, by the discrete
Stirling difference (or directly by factorial ratios). The total
\(R\)-reaction rate is \(O(\beta_n)\). Hence

\[
 \mathcal L_R(F_\ell-U_\theta)=o(\beta_n),          \tag{4.16}
\]

which proves (4.13). Finally, every bounded reaction has
\(|\Delta F_\ell|\le C(1+\log R_n)\), and every \(R\)-source propensity
is at most a fixed multiple of \(\beta_n\) on the tier. The reaction
menu is finite, proving (4.14). \(\square\)

\subsubsection{Proposition 4.3 (powered all-active
drift)}

Along every selected failed all-active sequence,

\[
 \mathcal L(1+F_\ell)^4(x_n)\longrightarrow-\infty. \tag{4.17}
\]

The conclusion is uniform over the finite failed-cone menu of a fixed
pair in the finite-exception sense.

\paragraph{Proof}

Absorb the \(C\beta_n\) term of (4.4) into (4.13). After decreasing
\(a_n\), still with \(a_n\to\infty\), Lemmas 4.1--4.2 give

\[
 \mathcal LF_\ell
 \le-\mathcal D_n-\beta_na_n                         \tag{4.18}
\]

and, for \(k=2,3,4\),

\[
 M_{k,n}:=\sum_r\lambda_r|\Delta_rF_\ell|^k
 \le C\{\mathcal D_n+\beta_n(1+\log^kR_n)\}.       \tag{4.19}
\]

The exact discrete binomial, or carré/Taylor, identity is

\[
\begin{split}
 \mathcal LW={}&4G^3\mathcal LF_\ell
       +6G^2\sum_r\lambda_r(\Delta_rF_\ell)^2\\
 &+4G\sum_r\lambda_r(\Delta_rF_\ell)^3
       +\sum_r\lambda_r(\Delta_rF_\ell)^4.
\end{split}\tag{4.20}
\]

The third moment is signed, so bounding it by its absolute moment can
only increase the right side. Substituting (4.18)--(4.19) gives

\[
\begin{split}
 \mathcal LW\le{}&-4G^3(\mathcal D_n+\beta_na_n)\\
 &+C(G^2+G+1)\mathcal D_n\\
 &+C\beta_n\{G^2(1+\log^2R_n)
      +G(1+\log^3R_n)+1+\log^4R_n\}.
\end{split}\tag{4.21}
\]

Because \(G\to\infty\), the second line is absorbed by
\(-G^3\mathcal D_n\). By (2.6), the ratio of the largest term in the
last line to \(G^3\beta_na_n\) is

\[
 {C(1+\log^2R_n)\over G a_n}\longrightarrow0,      \tag{4.22}
\]

and the other two ratios are smaller. Thus eventually

\[
 \mathcal LW(x_n)
 \le-2G_n^3(\mathcal D_n+\beta_na_n)\to-\infty.     \tag{4.23}
\]

This controls both the reversible top carré term near detailed balance
and the Taylor remainder away from it. It does not infer
convex-transform drift merely from \(\mathcal LF_\ell<0\).

Equivalently, if the harmless finite shifts in (4.8) are temporarily
deleted and \(q=\Delta_{yz}F_\ell=-h\), the two top directions pair as

\[
\begin{split}
 &(G+q)^4-G^4+e^q\{(G-q)^4-G^4\}\\
 &\quad=4G^3q(1-e^q)+6G^2q^2(1+e^q)
       +4Gq^3(1-e^q)+q^4(1+e^q).
\end{split}\tag{4.24}
\]

On the bounded flat-top ratio range, \(q(1-e^q)\le-cq^2\), while the
cubic term is nonpositive; hence this core is nonpositive for large
\(G\). Restoring the shifts costs \(O((A+B)s_nG^3)=O(\beta_nG^3)\) by
(4.11), which is absorbed by the \(-\beta_na_nG^3\) lower exit. This
paired form is the same domination as (4.18)--(4.23), and records
explicitly that the two enabled top directions are paired before
estimating their discrete curvature.

There are finitely many failed all-active descriptors for a fixed pair.
If no common finite exception existed, a violating divergent sequence
would have a subsequence in one descriptor and contradict (4.23). This
is the claimed finite-exception uniformity. \(\square\)

\subsection{\texorpdfstring{5. Fixed-class common-\(W\)
composition}{5. Fixed-class common-W composition}}

\subsubsection{Theorem 5.1 (audited prospective 26-pair
branch)}

For every selected support pair, every strongly connected orientation,
every positive rate vector, and every closed irreducible population
class \(\Gamma\), the physical mass-action CTMC is positive recurrent.

\paragraph{Proof}

Use the one function \(W\) in (2.5) on the whole class. Take any
divergent sequence in \(\Gamma\) and extract an exact source/D-tier
subsequence. An affine-infeasible descriptor cannot occur in the fixed
affine class. Every affine-feasible descriptor is therefore one of the
following.

\begin{enumerate}
\def\labelenumi{\arabic{enumi}.}
\tightlist
\item
  It passes the ordinary top-S/top-D criterion. The quantitative
  Anderson--Kim factorial estimate, with an arbitrary fixed linear
  correction, and the standard fourth-power remainder calculation give
  \(\mathcal LW\to-\infty\).
\item
  It is one of the 94 all-active failures. Proposition 4.3 gives
  \(\mathcal LW\to-\infty\).
\item
  It is one of the 30 one-active failures. There is no two-active failed
  descriptor by (2.1).
\end{enumerate}

The usual bad-sequence contradiction therefore gives a finite set
\(K_0\) such that \(\mathcal LW\le-1\) outside \(K_0\) and finitely many
fixed-width one-active tubes. A moving-boundary exit from such a tube
has two active coordinates and hence is in the generator-good region; no
promotion kernel or return to the old axis is needed.

On the all-species reflected marked chain update

\[
 D_i^+=(D_i+\zeta_i)^+,
 \qquad H_i=X_i-D_i.                                \tag{5.1}
\]

Then \(0\le D_i\le X_i\) and \(H_i\) is pathwise bounded above by its
initial value. In a divergent one-active bad tube, \(D_X=0\) would imply
\(X=H_X\) bounded, a contradiction. Thus every reachable divergent bad
tube has positive old debt. The two no-history rows of Lemma 3.1 are
incompatible with that marked history. Lemma 3.2 and the repeated-base
fourth-power theorem give, outside another finite set, an all-reaction
physical stopping time \(\tau>0\) with

\[
 \mathbb E_{x,d}[W(X_\tau)-W(x)+\tau]\le-1,         \tag{5.2}
\]

whose successful endpoint strictly reduces old \(D_X\). It uses the same
\(\ell\) and the same \(W\) as the generator-good region.

Generator steps and the episodes (5.2) are exactly the hypotheses of the
existing all-debt finite-target gluing theorem. Alternating them and
localizing at finite \(W\)-sublevels gives a nonnegative
supermartingale; bounded-index telescoping and Fatou yield a finite mean
hit of the finite marked target where every debt is zero. Its physical
projection is finite. The finite trace has a positive recurrent state,
and irreducibility of \(\Gamma\) promotes positive recurrence to every
state in the class.

Finally the process is nonexplosive. In a binary network, a reaction
with a bimolecular source cannot increase total population because
targets have molecularity at most two. Every population-increasing
propensity is therefore constant or linear. Stopping at population
levels and applying Gronwall controls the total population, while every
finite population sublevel has finitely many states and bounded total
jump rate. \(\square\)

\subsection{6. Reproduction and frozen
hashes}

The finite certificate freezes the following hashes:

\begin{CodeBlock}
one-active rows:       091dc59dc83b6445d36135a2d6b1e6c1bfb9565db88328fc1725c39f78bbc7e4
normalized profiles:  36b6eb1c1beadc939681fc7515e2adf3dfb9e16d116da4b8d2bd73ac77aabe12
all-active rows:       55669e3dc9b3fbb7ee84895ea158b3b596d339265cde1a23189bdf312a314bea
pair-fixed tops:       d1de16c1c45eaab3bce35a17db75bd947a8398a049722f100d96dd227b7e21d5
certificate payload:  a1f528caffb63729d889e792ce8831b899865e70c1564b2a0a33d46f295a39da
prospective remainder: 4e645eca1ba23849680f2f983e1fb8c9465001c5b0e8c0090a31a339bf18ec06
\end{CodeBlock}

Run

\begin{CodeBlock}
PYTHONPATH=src python3 -B src/prospective_26_candidate_theorem.py
PYTHONPATH=src python3 -B -m unittest \
  tests/test_prospective_26_candidate_theorem.py -v
\end{CodeBlock}

The machine checks exact supports, caps, categories, top reaction
vectors, cofactor tiers, and the discrete endpoint factorial identity.
It does not replace the analytic estimates above. The independent audit
replayed the one-active arbitrary-orientation scope extension, the
discrete powered-top remainder, the common-\(W\) gluing, nonexplosion,
and the exact zero-overlap arithmetic, and returned \textbf{PASS}.
Accordingly the scoped flags are

\begin{CodeBlock}
analytic_one_active_scope_extension_certified = true
analytic_powered_all_active_lift_certified = true
candidate_26_pair_recurrence_certified = true
global_t3_2_certified = false
\end{CodeBlock}
